


\newpage

%\footline={\hss\folio\hss}
%\pageno=1

\dspace
\subhead
4. Le groupe fondamental unipotent motivique
d'une vari\'et\'e rationnelle
\endsubhead

\subhead
4.1
\endsubhead
Soient $k$ un corps de nombres, $P$ une droite
projective sur $k$, $X$ le compl\'ement d'un
ensemble fini $S$ de points rationnels et $x,y\in
X(k)$.
Dans Deligne (1989) \S{13}, nous avons d\'efini un
``espace de (classes d'homotopie de) 
chemins motivique'' not\'e $P_{y,x}$.
Sous les hypoth\`eses ci-dessus, et apr\`es
oubli de l'aspect cristallin de loc. cit.,
c'est un sch\'ema affine en la cat\'egorie
tannakienne $\scrR^{H+\ell}$ de 2.15.
Nous le noterons ici $P_{y,x}^{H+\ell}$, ou
$P_{y,x}^{H+\ell}(X)$ s'il y a lieu de
pr\'eciser $X$.
On dispose d'une ``composition des chemins''
$P_{z,y}^{H+\ell}\times P_{y,x}^{H+\ell}
\to P_{z,x}^{H+\ell}$.
Posons $\pi_1^{H+\ell}(X,x):=P_{x,x}^{H+\ell}$.
Pour la composition des chemins, c'est un
$\scrR^{H+\ell}$-sch\'ema en groupes affine et
$P_{y,x}^{H+\ell}$ est un espace principal
homog\`ene \`a droite sous
$\pi_1^{H+\ell}(X,x)$, \`a gauche sous
$\pi_1^{H+\ell}(X,y)$: un
$(\pi_1^{H+\ell}(X,x),
\pi_1^{H+\ell}(X,y))$-bitorseur.
Pour $\sigma$ un plongement de $k$ dans une
cl\^{o}ture alg\'ebrique $C$ de $\dbR$, la
r\'ealisation $\sigma$ de
$\pi_1^{H+\ell}(X,x)$ est l'enveloppe
alg\'ebrique prounipotente de $\pi_1(X(C),x)$.
Si $I$ est l'id\'eal d'augmentation de
l'alg\`ebre de groupe $\dbQ[\pi_1(X(C),x)]$, c'est
le spectre de l'alg\`ebre de Hopf commutative
$$
\colim(\dbQ[\pi_1(X(C),x)]/I^N)\upvee.
$$
i.e. c'est l'enveloppe alg\'ebrique
pro-unipotente de $\pi_1(X(C),x)$.

\subhead
4.2
\endsubhead
Soit $A_{y,x}^{H+\ell}$ l'alg\`ebre affine
(2.6) de $P_{y,x}^{H+\ell}$.
D'apr\`es 3.12, 3.13, le 
Ind-objet $A_{y,x}^{H+\ell}$ de
$\scrR^{H+\ell}$ est l'image par le foncteur
r\'ealisation d'un Ind-objet $A_{y,x}$ de
$MT(k)$.
D'apr\`es 2.14, 2.15 (ii), le produit
$A_{y,x}^{H+\ell}\otimes A_{y,x}^{H+\ell}\to
A_{y,x}^{H+\ell}$ provient d'un produit sur
$A_{y,x}$, et $P_{y,x}^{H+\ell}$ est la
r\'ealisation du $MT(k)$-sch\'ema affine
$P_{y,x}:=\Spec(A_{y,x})$.
De m\^{e}me, la composition des chemins est
d\'efinie par un coproduit $A_{z,x}^{H+\ell}\to
A_{z,y}^{H+\ell}\otimes A_{y,x}^{H+\ell}$, ce
coproduit se rel\`eve \`a $MT(k)$ et d\'efinit
un morphisme $P_{z,y}\times P_{y,x}\to
P_{z,x}$.

Nous appellerons ici $P_{y,x}$ l'{\it espace
de chemins de $x$ \`a $y$ motivique}.
Nous appelerons $P_{x,x}$ le groupe
fondamental motivique de $X$ en $x$ et le
noterons $\pi_1^{\mot}(X,x)$ ou simplement
$\pi_1(X,x)$.
Mise en garde: bien que la notation ne
l'indique pas, il s'agit de $\pi_1$ rendus
unipotents.

\subhead
4.3
\endsubhead
Appelons ``point base de $X$'' soit un point
rationnel de $X=P-S$, soit un vecteur tangent
non nul en un point de $S$ (``points base \`a
l'infini'').
Pour $x$ et $y$ deux points-base de $X$, 
Deligne (1989) d\'efinit \S{15} un
``espace de chemins motivique de $x$ \`a
$y$''.
Comme en 4.1, c'est, apr\`es oubli d'un
aspect cristallin, un
$\scrR^{H+\ell}$-sch\'ema affine.
Nous le noterons ici $P_{y,x}^{H+\ell}$, et,
pour $x=y$, $\pi_1^{H+\ell}(X,x)$.

\proclaim{Th\'eor\`eme 4.4}
Avec les notations de 4.1 et 4.3, si $x$ et
$y$ sont deux points-base de $X$, le
$\scrR^{H+\ell}$-sch\'ema $P_{y,x}^{H+\ell}$
est motivique, i.e. la r\'ealisation d'un
$MT(k)$-sch\'ema $P_{y,x}$.
\endproclaim

Si $S$ est vide ou r\'eduit \`a un point,
$X=P-S$ est simplement connexe, les
$P_{y,x}^{H+\ell}$ sont r\'eduits \`a un point
et 4.4 est trivial.
Nous pouvons donc supposer,
et supposerons, que $\vert S\vert\Ge
2$.
Le cas o\`u $x$ et
$y$ sont \`a distance finie ayant d\'ej\`a
\'et\'e trait\'e, nous supposerons
au moins l'un de $x$ et $y$ \`a l'infini.

\medskip\noindent
{\sans Cas 1.}
$x\in X(k)$ et $y$ est un vecteur tangent non
nul en $\ybar\in S$.

Soit $Z$ le $\MT(k)$-sch\'ema
des homomorphismes de $\dbQ(1)$
dans $\pi_1^{\mot}(X,x)$.
Identifiant homomorphismes de sch\'emas en
groupe unipotents, et homomorphismes entre
leurs alg\`ebres de Lie, on peut en donner la
description suivante: Lie $\pi_1^{\mot}(X,x)$
est une pro-alg\`ebre de Lie dans $MT(k)$, en
particulier un pro-objet de $MT(k)$ (c'est la
limite projective des Lie~$(\pi_1^{\mot}(X,x)/\grZ^N)$, pour $\grZ$ la
s\'erie centrale descendante),
et on d\'efinit $Z$ comme \'etant le
$MT(k)$-sch\'ema pro-vectoriel (2.6) Lie
$\pi_1^{\mot}(X,x)(-1)$.

Dans $Z^{H+\ell}:=\real^{H+\ell}(Z)$,
on dispose du sous-sch\'ema ferm\'e ``classe de
conjugaison de la monodromie locale autour de
$\ybar\,$''.
Voici sa d\'efinition.
Ecrivons $\dbQ(1)^{H+\ell}$ pour
$\real^{H+\ell}(\dbQ(1))$.
On dispose d'un morphisme ``monodromie locale
autour de $\ybar\,$'':
$$
\dbQ(1)^{H+\ell}\to\pi_1^{H+\ell}(X,y).
\tag{4.4.1}
$$
Dans (4.4.1), $\dbQ(1)^{H+\ell}$ est vu comme
un $\scrR^{H+\ell}$-sch\'ema vectoriel, avec sa
structure de groupe.
Si $T_{\ybar}$ est l'espace tangent en
$\ybar$, et que
$T_{\ybar}^*:=T_{\ybar}-\{0\}$, il y a lieu de
voir (4.4.1) comme un morphisme
$$
\pi_1^{H+\ell}(T_{\ybar}^*,y)\to\pi_1^{H+\ell}
(X,y).
$$
La donn\'ee de (4.4.1) \'equivaut \`a celle de
$$
\dbQ(1)^{H+\ell}\to\Lie\,\pi_1^{H+\ell}
(X,y).
\tag4.4.2
$$
Dans (4.4.2), $\dbQ(1)^{H+\ell}$ est un objet
de $\scrR^{H+\ell}$, muni de sa structure
d'alg\`ebre de Lie commutative.

On dispose aussi de l'espace de chemins
$P_{y,x}^{H+\ell}$ et, par composition des
chemins, de
$$
(p,\gamma)\longmapsto p^{-1}\gamma p\colon\,
P_{y,x}^{H+\ell}\times\pi_1^{H+\ell}(X,y)\to
\pi_1^{H+\ell}(X,x).
$$
Composant avec (4.4.1), on en d\'eduit
$$
P_{y,x}^{H+\ell}\times\dbQ(1)^{H+\ell}\to
\pi_1^{H+\ell}(X,x).
\tag4.4.2
$$

Le groupe $\pi_1^{H+\ell}(X,y)$ agit par
composition \`a gauche sur
$P_{y,x}^{H+\ell}$.
Divisant \`a gauche par
$\dbQ(1)^{H+\ell}$, vu par (4.4.1) comme
sous-groupe de $\pi_1^{H+\ell}(X,y)$, on d\'efinit
$$
P_{\ybar,x}^{H+\ell}:=\dbQ(1)^{H+\ell}\setminus
P_{y,x}^{H+\ell}
$$
et (4.4.2) se factorise par
$P_{\ybar,x}^{H+\ell}\times\dbQ(1)^{H+\ell}$,
d\'efinissant
$$
P_{\ybar,x}^{H+\ell}\to Z^{H+\ell}\tag4.4.3
$$

\proclaim{Lemme 4.5}
Le morphisme (4.4.3) est un plongement ferm\'e
de $\scrR^{H+\ell}$-sch\'emas.
\endproclaim

\demo{Preuve}
Il suffit de le v\'erifier en r\'ealisation de
de Rham, o\`u l'assertion se r\'eduit \`a la
suivante.
\enddemo

\proclaim{Lemme 4.6}
Soit $\scrL$ l'alg\`ebre de Lie sur $k$
engendr\'ee  par des \'el\'ements $(e_s)_{s\in
S}$ soumis \`a la seule relation que $\sum
e_s=0$.
On suppose $\vert S\vert\Ge 2$, on fixe
$\ybar\in S$ et on pose $e:=e_{\ybar}$.
Soit $G$ le sch\'ema en groupe prounipotent
limite projective des
$G_N:=\exp(\scrL/\grZ^N)$.
Le morphisme de $G/\exp(ke)$ dans $\Lie(G)$:
$g\mapsto \ad\,g(e)$, identifie $G/\exp(ke)$
\`a un sous-sch\'ema ferm\'e de $\Lie(G)$,
l'orbite adjointe de $e$.
\endproclaim

Dans 4.6, $\Lie(G)$ est le sch\'ema pro-vectoriel
limite projective des sch\'emas vectoriels
$\Lie(G_N)$.

\demo{Preuve}
Les orbites de l'action d'un groupe unipotent
sur un sch\'ema affine sont ferm\'ees Demazure,
Gabriel (1970) IV 2.2.7).
Si $e_N$ est l'image de $e$ dans
$\scrL/\grZ^N$ et $Z_N\Subset G_N$ son
fixateur, l'application $g\mapsto \ad\,g(e_N)$
de $G_N/Z_N$ dans $\Lie(G_N)$ est donc un
plongement ferm\'e.
Par passage \`a la limite projective,
l'application $g\mapsto\ad\,g(e)$ de
$\lim\,G_N/Z_N$ dans $\Lie(G)=\lim\,\Lie(G_N)$
est \'egalement un plongement ferm\'e.

L'alg\`ebre de Lie du 
sous-groupe $Z_N$ de $G_N$ est le
centralisateur de $e_N$ dans
$\scrL/\grZ^N$.
Fixons $b$ dans $S$ distinct de $\ybar$.
L'alg\`ebre de Lie $\scrL$ est librement
engendr\'ee par les $e_s$ ($s\not=b$).
D'apr\`es 4.7 ci-dessous, le centralisateur de
$e$ est r\'eduit \`a $ke$.
Munissons $\scrL$ de la graduation pour
laquelle les $e_s$ sont de degr\'e $1$.
La sous-alg\`ebre $\grZ^N$ est la
sous-alg\`ebre ``degr\'e $\Ge N$''.
L'application lin\'eaire
$\ad\,e\colon\,\scrL\to\scrL$ \'etant de
degr\'e $1$ et de noyau r\'eduit \`a $ke$, on a
$$
\Lie\,Z_N=\Ker(\ad\,e_N\colon\,\scrL/\grZ^N\to
\scrL/\grZ^N)=ke_N+\grZ^{N-1}/\grZ^N.
$$

La projection de $G_{N+1}/Z_{N+1}$ sur 
$G_N/Z_N$ se factorise donc par
$G_N/\exp(ke_N)$ et
$$
G/\exp(ke)=
\lim\,G_N/\exp(ke_N)\simover \lim\,G_N/Z_N(e).
$$
\enddemo

\proclaim{Lemme 4.7}
Dans une alg\`ebre de Lie libre $\scrL$ sur
$k$, le centralisateur d'un g\'en\'erateur $e$
est r\'eduit aux multiples de ce
g\'en\'erateur.
\endproclaim

D'apr\`es Reutenauer (1993) 2.10, 
dans une alg\`ebre de Lie
libre, deux \'el\'ements non nuls quelconques
qui commutent sont multiples l'un de l'autre.
Le cas particulier 4.7 est, plus simplement,
cons\'equence de ce que l'alg\`ebre de Lie
libre $\Lib(\{e\}\CUP A)$ est le produit
semi-direct de $ke$ et de
$\Lib((\ad\,e^n(a))_{a\in A,n\Ge 0}$ et de ce
que la d\'erivation $\ad\,e^n(a)\mapsto
\ad\,e^{n+1}(a)$ de l'alg\`ebre associative
libre engendr\'ee par les $\ad\,e^n(a)$ est
injective.

Il r\'esulte de 4.5 que $P_{\ybar,x}^{H+\ell}$
est motivique.
Soient en effet $A$ l'alg\`ebre de $Z$ et
$A^{H+\ell}:=\real^{H+\ell}(A)$ celle de
$Z^{H+\ell}$.
D'apr\`es 4.5, $P_{\ybar,x}^{H+\ell}$ est
isomorphe \`a un sous-sch\'ema ferm\'e de
$Z^{H+\ell}$, d\'efini par un id\'eal
$a^{H+\ell}$ de $A^{H+\ell}$.
D'apr\`es  2.14, 2.15, cet id\'eal est la
r\'ealisation d'un id\'eal $a$ de $A$, et
$P_{\ybar,x}^{H+\ell}$ est isomorphe \`a la
r\'ealisation de $\Spec(A/a)$.


\subhead
4.8 Fin du cas $1$ de la preuve de 4.4
\endsubhead
Soit $b\in S$ distinct de $\ybar$, et
$X':=P-\{\ybar,b\}$.
Identifions $\dbP$ et $\dbP^1$, par un
isomorphisme envoyant $\ybar$ sur $0$ et $b$
sur $\infty$, donc $X'$ sur
$A^*=\Spec\,k[u,u^{-1}]$.
Soit $t:=$ du $(y)$ la coordonn\'ee du vecteur
tangent $y$ en $u=0$, et notons encore $t$ le
point de $\dbG_m$ de coordonn\'ee $t$.
Les espaces de chemins $P_{y,x}^{H+\ell}(X')$
et $P_{t,x}^{H+\ell}(\dbG_m)$ sont
canoniquement isomorphes (Deligne (1989) \S{14}).
Puisque $t$ et $x$ sont des points-base \`a
distance finie de $\dbG_m$, $P_{y,x}^{H+\ell}$
est motivique.
C'est d'ailleurs simplement un torseur de
Kummer: on a $\pi_1^{\mot}(\dbG_m,x)=\dbQ(1)$,
$t/x\in k^*$ d\'efinit une extension de
$\dbQ(0)$ par $\dbQ(1)$, donc un
$\dbQ(1)$-torseur, et $P_{y,x}^{H+\ell}$ 
est la r\'ealisation de ce torseur.

\proclaim{Lemme 4.9}
Le morphisme
$$
P_{y,x}^{H+\ell}(X)\to
P_{y,x}^{H+\ell}(X')\times
P_{\ybar,x}^{H+\ell}(X)
$$
$(1^{\text{\rm en}}$ facteur: fonctorialit\'e
en $X$; $2^{\text{\rm e}}$ facteur: passage ou
quotient$)$ est un isomorphisme.
\endproclaim

Puisque tant $P_{y,x}^{H+\ell}(X')$ que
$P_{\ybar,x}^{H+\ell}(X)$ sont motiviques 4.9 
implique que $P_{y,x}^{H+\ell}(X)$ est motivique.

\demo{Preuve de 4.9}
Il suffit de v\'erifier 4.9 en r\'ealisation
de de Rham.
Avec les notations de 3.6, l'assertion se
r\'eduit \`a la suivante.
\enddemo

\proclaim{Lemme 4.10}
Regardons $\dbG_a=\exp(ke)$ comme un quotient
de $G$ par $e\mapsto e$, $e_a\mapsto 0$ pour
$a\not=\ybar,b$.
L'application
$$
G\to \exp(ke)\times\exp(ke)\setminus G
$$
est un isomorphisme.
\endproclaim

\demo{Preuve}
$G$ est le produit semi-direct de $\exp(ke)$
par le noyau de la projection sur $\exp(ke)$,
cf. la preuve de 4.7.
\enddemo

\subhead
4.11
\endsubhead
Pour finir la preuve de 4.4 il reste, compte
tenu de ce que $P_{y,x}^{H+\ell}$, est
isomorphe comme $\scrR^{H+\ell}$-sch\'ema \`a
$P_{x,y}^{H+\ell}$, \`a traiter le

\medskip\noindent
{\sans Cas 2.}
$x$ et $y$ sont des points base \`a l'infini.

Choisissons un point base $z\in X(k)$.
On a une expression de $P_{y,x}^{H+\ell}$
comme compos\'e de bitorseurs
$$
P_{y,x}^{H+\ell}=P_{y,z}^{H+\ell}
\times_{\pi_1^{H+\ell}(X,z)}  P_{z,x}^{H+\ell}
$$
o\`u au second membre chaque
$\scrR^{H+\ell}$-sch\'ema est motivique.
Que $P_{y,x}^{H+\ell}$ soit motivique en
r\'esulte.

\subhead
4.12
\endsubhead
Soit $X$ une vari\'et\'e alg\'ebrique lisse
sur un corps de nombres $k$, et $x,y\in X(k)$.
On suppose $X$ s\'epar\'e.
Sous l'hypoth\`ese que $H^1(\Xbar,\scrO)=0$
pour $\Xbar$ une compactification lisse de $X$,
nous avons d\'efini dans Deligne (1989) \S{13} un
``espace de chemins'' $P_{y,x}$ qui est un
sch\'ema en une cat\'egorie de syst\`emes de
r\'ealisations.
Nous le noterons ici,
apr\`es omission de l'aspect cristallin,
$P_{y,x}^{\real}$.
Comme pr\'ec\'edemment, les $P_{y,x}^{\real}$
forment un groupo{\umi}de.

\proclaim{Th\'eor\`eme 4.13}
Si la vari\'et\'e $X$ est unirationnelle,
$P_{y,x}^{\real}$ est motivique,
r\'ealisation d'un $\MAT(k)$-sch\'ema.
\endproclaim

Les cat\'egories de syst\`emes de
r\'ealisations consid\'er\'es sont des champs
sur $\Spec(k)_{\et}$.
Le th\'eor\`eme \'equivaut donc \`a dire qu'il
existe une extension galoisienne $k_1$ de $k$
telle qu'apr\`es extension du corps de base
\`a $k'$ dans $\scrC(k_1/k)$,
$P_{y,x}^{\real}(X_{k'})$ est motivique,
r\'ealisation d'un $\MT(k')$-sch\'emas.

\demo{\bf Preuve pour $X$ un ouvert de $\dbP^n$}
Soit $D\Subset\dbP^n$ une droite, assez
g\'en\'erale pour que le morphisme de
sch\'emas en groupe unipotents
$$
\pi_1^{\real}(X\CAP D,x)\to
\pi_1^{\real}(X,x)
$$
($x$ quelconque dans $X\CAP D$) soit surjectif.
\enddemo

Si $X\CAP D$ est le compl\'ement dans $D$ d'un
ensemble fini de points rationnels,\break
$\pi_1^{\real}
(X\CAP D,x)$ est dans $\scrR^{H+\ell}$ et,
d'apr\`es 4.4, motivique, r\'ealisation
de $\pi_1^{\mot}(X\CAP D,x)$ dans $MT(k)$.
L'alg\`ebre affine de $\pi_1^{\real}(X,x)$ est
une sous-alg\`ebre de celle de $\pi_1^{\real}(X
\CAP D,x)$, donc est \'egalement motivique,
dans $MT(k)$.
En g\'en\'eral, pour $k'$ assez grand (i.e.,
dans un $\scrC(k_1/k)$ convenable), $X\CAP D$
deviendra apr\`es extension du corps de base
de $k$ \`a $k'$ le compl\'ement d'un ensemble de
points rationnels sur $k'$, et la construction
ci-dessus nous fournit
$\pi_1^{\mot}(X_{k'},x)$ dans $MT(k')$,
fonctoriel en $k'$, i.e. $\pi_1^{\mot}(X,x)$
dans $\MAT(k)$.

Soit $y\in X(k)$, distinct de $x$.
Prouvons que $P_{y,x}^{\real}$ est motivique.
Soit $D$ la droite de $\dbP^n$ passant par $x$
et $y$.
Comme pr\'ec\'edemment, on se ram\`ene \`a
supposer que $X\CAP D$ est le compl\'ement, dans
$D$, d'un ensemble fini de points rationnels.
Le $\scrR^{H+\ell}$-sch\'ema
$P_{y,x}^{\real}(X\CAP D)$ est alors la
r\'ealisation de $P_{y,x}^{\mot}(X\CAP D)$
dans $\MT(k)$.
On obtient $P_{y,x}^{\mot}(X)$ en poussant le
$\pi_1(D\CAP X,x)$-torseur
$P_{y,x}^{\mot}(D\CAP X)$ par
$$
\pi_1(D\CAP X,x)\to\pi_1(X,x).
$$

Pour $y,z\in X(k)$, $P_{y,z}^{\mot}$ est
d\'eduit
des $\pi_1^{\mot}(X,x)$-torseurs
$P_{y,x}^{\mot}$ et $P_{z,x}^{\mot}$.

\demo{\bf Preuve de $4.13$ {\rm (cas g\'en\'eral)}}
La vari\'et\'e $X$, \'etant unirationelle, est
domin\'ee par un ouvert $U$ de $\dbP^n$.
Fixons $x\in U(k)$.
Pour $\sigma$ un plongement complexe de $k$,
le morphisme
$\pi_1(U(\dbC),x)\to\pi_1(X(\dbC),x)$ a une
image d'indice fini.
Le morphisme induit sur les enveloppes
alg\'ebriques unipotentes est donc surjectif,
et l'alg\`ebre affine de $\pi_1^{\real}(X,x)$
est motivique en tout que sous-alg\`ebre de
l'alg\`ebre affine de $\pi_1^{\real}(X,x)$,
que nous savons d\'ej\`a \^{e}tre motivique.
Le cas des $P_{y,x}$ se traite ensuite comme
plus haut, en utilisant que deux points de $X$
peuvent \^{e}tre connect\'es par une 
cha\^{\dotless\char'020}ne de courbes
rationnelles.
\enddemo

\remark{\bf Remarque 4.14}
Utilisant 4.4, on pourrait aussi dans 4.13
prendre pour $x$ et $y$ des points-base \`a
l'infini, au sens de Deligne (1989) \S{15}.
\endremark

\proclaim{Proposition 4.15}
Supposons que $\Xbar$ soit une
compactification normale de la vari\'et\'e
unirationelle $X$, et soient $x,y\in X$.
Si les composantes irr\'eductibles de
codimension un dans $X$ de $Y:=\Xbar-X$ sont
absolument irr\'eductibles, alors
$P_{y,x}^{\mot}$ est un $\MT(k)$-sch\'ema.
\endproclaim

\demo{Preuve}
Supposons d'abord que $x=y$.
Pour que $\pi_1(X,x)$ soit un
$\MT(k)$-sch\'ema il suffit que
$\Lie\,\pi_1(X,x)$ soit dans $\MT(k)$.
D'apr\`es 2.18, il suffit m\^{e}me que son
gradu\'e par le poids le soit.
Ce gradu\'e \'etant engendr\'e par $H_1(X)$,
purement de poids $-2$, il suffit que l'action
de $\Gal(\kbar/k)$ sur $\omega H_1(X)$ soit
triviale.
Ce groupe est
$$
\scrO^*(X_{\kbar})/\kbar^*\otimes\dbQ
$$
et ses \'el\'ements sont d\'etermin\'es par
leurs valuations le long des composantes
irr\'eductibles de codimension un de
$Y_{\kbar}$.
Si le groupe $\Gal(\kbar/k)$ ne les permute
pas, il agit donc trivialement sur
$\scrO^*(X_{\kbar})/\kbar^*\otimes\dbQ$.
\enddemo

Le cas g\'en\'eral r\'esulte du

\proclaim{Lemme 4.16}
Si un $\MAT(k)$-sch\'ema $P$ est un torseur
sous un $\MT(k)$-sch\'ema en groupes unipotent
$G$, il est lui-m\^{e}me un $\MT(k)$-sch\'ema.
\endproclaim

\demo{Preuve}
Le groupe $G$ \'etant unipotent, si on le fait
agir sur son alg\`ebre affine $A$ par
translations \`a droite, celle-ci admet une
filtrations exhaustive Fil par des sous-motifs
stables sous $G$, et tels que $G$ agisse
trivialement sur le gradu\'e associ\'e.
L'alg\`ebre affine $A_P$ de $P$, d\'eduite de
$A$ en poussant par le $G$-torseur $P$, admet
donc une filtration, encore not\'ee Fil, telle
que $\Gr^{\Fil}(A)\simover \Gr^{\Fil}(A_P)$.
Il reste \`a appliquer 2.18.
\enddemo

\proclaim{Proposition 4.17}
Soit $S$ un ensemble fini de places finies de
$k$ et supposons que la vari\'et\'e
unirationelle $X$ est la fibre g\'en\'erale du
compl\'ement $X_{\scrO}$, dans $\Xbar_{\scrO}$
propre et lisse sur $\Spec(\scrO_S)$, d'un
diviseur \`a croisements normaux relatif,
r\'eunion de diviseurs lisses, chacun de fibre
g\'en\'erale absolument irr\'eductible.
Si $x$ et $y$ dans $X(k)$ proviennent de
$x_{\scrO}$, $y_{\scrO}$ dans
$X_{\scrO}(\scrO_S)$, alors $P_{y,x}^{\mot}$
est dans $\MT(\scrO_S)$.
\endproclaim

\demo{Preuve}
Pour tout nombre premier $\ell$, l'hypoth\`ese
assure que la r\'ealisation $\ell$-adique de
$P_{y,x}^{\mot}$ est non ramifi\'ee en dehors
de $S$ et des places au-dessus de $\ell$, et on
applique 1.5.
\enddemo

\remark{\bf 4.18 Remarque}
Un \'enonc\'e analogue vaut pour $x$ ou $y$ \`a
l'infini.
\endremark


\subhead
4.19
\endsubhead
Soit $X$ une vari\'et\'e lisse unirationnelle
sur un corps de nombres $k$.
Le th\'eor\`eme 4.13 permet de d\'efini la
cat\'egorie 
$\MAT(X/k)$ des {\it syst\`emes locaux
unipotents de motifs d'Artin-Tate mixtes} sur
$X$.
On choisit $x\in X(k)$ et on d\'efinit un
objet $M$ de $\MAT(X/k)$ comme \'etant la
donn\'ee de $M_x$ dans $\MAT(k)$ muni d'une
action de $\pi_1(X,x)$.
Il y a lieu de voir $M_x$ comme \'etant la
fibre de $M$ en $x$.
Le choix de $x$ n'importe pas: pour $y$ dans
$X(k)$, on d\'efinit $M_y$ comme \'etant le
tordu de $M_x$ par le $\pi_1(X,x)$-torseur
$P_{y,x}$.
Il reviendrait au m\^{e}me de d\'efinir
$\MAT(X/k)$ comme la cat\'egorie des 
repr\'esentations, dans $\MAT(k)$, du
groupo{\umi}de des $P_{y,x}$ pour $x,y\in
X(k)$.
On pourrait aussi prendre des points base \`a
l'infini.

Puisque $\Lie\,\pi_1(X,x)$ est en poids $<0$,
$\pi_1(X,x)$ agit trivialement sur
$\Gr^W(M_x)$ et, pour $x$ et $y$ deux  points
base de $X$, $\Gr^W(M_x)$ et $\Gr^W(M_y)$ sont
canoniquement isomorphes.
Le foncteur fibre $\omega(M):=\omega(M_x)$ est
donc 
ind\'ependant de $x$.
On d\'efinit la sous-cat\'egorie $\MT(X)$ de
$\MAT(X/k)$ par la condition ``$M_x$ est dans
$\MT(k)$''.
Le choix de $x$ n'importe pas (2.18).

\subhead
4.20
\endsubhead
Soit $\sigma$ un plongement de $k$ dans une
cl\^{o}ture alg\'ebrique $C$ de $\dbR$.
Pour $M$ dans $\MAT(X/k)$, $(M_x)_\sigma$ est
une structure de Hodge-Tate mixte, munie d'une
action de Hodge mixte de $\pi_1(X,x)_\sigma$.
D'apr\`es Hain and Zucker (1987) (5.21),
ces donn\'ees d\'efinissent
une variation admissible $M_\sigma$ de
structures de Hodge mixtes sur $X(C)$.

Le torseur $(P_{y,x})_\sigma$ tord $(M_\sigma)_x$ 
en $(M_\sigma)_y$ et la variation $M_\sigma$ ne 
d\'epend donc pas du choix de $x$.

De m\^{e}me, $M$ d\'efinit un
$\dbQ_\ell$-faisceaux lisse sur $X$.

\subhead
4.21
\endsubhead
R\'eciproquement, supposons donn\'e sur $X$ un
syst\`eme de r\'ealisations $M$, candidat \`a
\^{e}tre la r\'ealisation d'un syst\`eme local
unipotent de motifs de Tate mixtes, de
gradu\'e par le poids constant.
Alors, si en un point $x$ de $X$ ce syst\`eme
est motivique, il est lui-m\^{e}me motivique.
En effet, dans les diverses r\'ealisations, il
d\'efinit une repr\'esentation de
$\pi_1(X,x)$.
Appliquant 2.14, on obtient une
repr\'esentation de $\pi_1^{\mot}(X,x)$ sur
le motif de r\'ealisation $M_x$, et cette
repr\'esentation d\'efinit le syst\`eme local
voulu.

\subhead
4.22
\endsubhead
Si $K$ est le corps des fonctions rationnelles
sur $X$, il est raisonnable de d\'efini
$\MT(K)$ comme la limite inductive des
cat\'egorie $\MT(U)$, pour $U$ un ouvert non
vide de plus en plus petit de $X$.
La relation entre cette cat\'egorie et la
cat\'egorie d\'eriv\'ee motivique $\DM(K)$ de
Voevodsky n'est pas claire.
Faute de disposer de la conjecture
d'annulation de Beilinson-Soul\'e pour $K$, on
ne sait pas extraire de $\DM(K)$ une
cat\'egorie ab\'elienne de motifs de Tate
mixtes sur $K$.
Pour $X$ rationnelle de dimension un, la
conjecture d'annulation est vraie, mais
nous n'avons pas prouv\'e l'\'equivalence des
deux constructions.
Pour la cat\'egorie $\MT(K)$ d\'efinie ici, on
a
$$
\alignat2
&\Ext^1(\dbQ(0),\dbQ(1))=K^*\otimes\dbQ
  &\qquad&\text{\rm et}\\
&\Ext^1(\dbQ(0),\dbQ(n))=K_{2n-1}(k) 
  &\qquad&\text{\rm pour $n\Ge 2$}.
\endalignat
$$

\subhead
4.23 Remarque
\endsubhead
Soit $\dbQbar$ une cl\^{o}ture alg\'ebrique de
$\dbQ$ et d\'efinissons $\MT(\dbQbar)$ 
comme la limite inductive des $\MT(F)$, pour
$F$ une extension finie de $\dbQ$ dans
$\dbQbar$.
Goncharov (1994) conjecture que, dans
cette limite, les $P_{y,x}(\dbP^1-S)$
engendrent tous les motifs de Tate
mixte.
En d'autres termes: que tout motif de Tate mixte
sur $\dbQbar$ est un sous-quotient d'un
produit tensoriel de sous-motifs d'alg\`ebres
affines $\scrO(P_{y,x}(\dbP^1-S))$ et de duaux
de tels sous-motifs.

Les remarques suivantes montrent qu'il est en
tout cas difficile de sortir de la
sous-cat\'egorie tannakienne $\scrT$ de
$\MT(\dbQbar)$ engendr\'ee par les
$P_{y,x}(\dbP^1-S)$.

\medskip\noindent
{\sans (a)}\enspace
Pour $X$ lisse unirationnelle, les
$P_{y,x}(X)$ sont dans $\scrT$.
Cela r\'esulte de la preuve de 4.13.

\smallskip\noindent
{\sans (b)}\enspace
Si $M$ est un syst\`eme local unipotent de
motifs de Tate mixte sur $X$ lisse
unirationnelle, et si en un point $M$ est dans
$\scrT$, alors $M$ est en chaque point dans
$\scrT$.
En effet, $M_y$ est tordu de $M_x$ par
$P_{y,x}$.
Il se plonge donc dans
$M_x\otimes\scrO(P_{y,x})$, et est dans
$\scrT$ puisque $M_x$ et $\scrO(P_{y,x})$ le
sont.


